We now prove a result which is sometimes called the \emph{analytic Fredholm theorem}. This is the basis for two important results: the Fredholm alternative theorem (Theorem~\ref{thm:fredholm-alternative-compact}) and the Hilbert-Schmidt theorem (Theorem~\ref{thm:hilbert-schmidt-spectral}).

\begin{Theorem}{Analytic Fredholm theorem}
\label{thm:analytic-fredholm}
Let $G\subset\mathbb{C}$ be a domain. Suppose the mapping
\begin{equation}\label{eq:analytic-compact-fredholm-operator-family}
\mathbb{C}\supset G\ni\lambda\mapsto B(\lambda)\in\mathcal{L}(H)
\end{equation}
is analytic on $G$ and that $B(\lambda)$ is compact at each $\lambda\in G$. Then, either\begin{enumerate}\item $(I-B(\lambda))^{-1}$ exists for no $\lambda\in G$, or\item $(I-B(\lambda))^{-1}$ exists for every $\lambda\in G\backslash S$ where $S$ is a discrete set in $G$ (i.e., it has no limit point in $G$). In this case the function $\lambda\mapsto(I-B(\lambda))^{-1}$ is analytic on $G\backslash S$, and if $\lambda\in S$, then $B(\lambda)\psi=\psi$ has a finite-dimensional family of solutions.\end{enumerate}
\end{Theorem}

\begingroup
\postdisplaypenalty=10000
\renewcommand{\endProof}{\par\nobreak\hfill$\blacksquare$\par\vspace{0.1cm}}
\begin{Proof}
We give the proof only in a neighborhood of a point $\lambda_0\in G$. Standard connectedness arguments can be used to extend the result to all of $G$.

Let $\lambda_0\in G$ be given. Since $\lambda\mapsto B(\lambda)$ is continuous, we can choose $r>0$ such that
\begin{equation}\label{eq:analytic-fredholm-local-continuity-bound}
\|B(\lambda)-B(\lambda_0)\|<\dfrac{1}{2}
\end{equation}
for all $\lambda$ in the disk $D_r=\{\lambda\in G\mid|\lambda-\lambda_0|<r\}$. Using the construction of Theorem~\ref{thm:compact-hilbert-finite-rank-approximation}, we see that there is an operator of finite rank $B_N$ such that
\begin{equation}\label{eq:analytic-fredholm-finite-rank-approximation-bound}
\|B_N-B(\lambda_0)\|<1/2.
\end{equation}
Now, using the geometric series techniques of the proof of Theorem~\ref{thm:bounded-operator-analytic-resolvent}, the reader can verify that
\begin{equation}\label{eq:analytic-fredholm-small-remainder-inverse}
(I-B(\lambda)+B_N)^{-1}
\end{equation}
exists as a bounded operator and is analytic on $D_r$.

Now let
\begin{equation}\label{eq:analytic-fredholm-finite-rank-reduction}
F(\lambda):=B_N\circ(I-B(\lambda)+B_N)^{-1}.
\end{equation}
Note that
\begin{equation}\label{eq:analytic-fredholm-operator-factorization}
(I-B(\lambda))=(I-F(\lambda))(I-B(\lambda)+B_N).
\end{equation}
Thus, $I-B(\lambda)$ is invertible if and only if $I-F(\lambda)$ is. However, $F$ has finite rank, so by Lemma~\ref{lem:analytic-finite-rank-alternative}, $I-F(\lambda)$ is either invertible at no $\lambda\in G$ or is invertible off a discrete $S\subset G$.

The proof that the solution space of $B(\lambda)\psi=\psi$ is finite-dimensional follows from the compactness of $B(\lambda)$ and is left to the reader (Problem~\ref{ex:compact-nonzero-eigenspace-finite-dimensional}). This completes the proof.
\end{Proof}
\endgroup

We now use the analytic Fredholm theorem to derive the following characterization of the spectrum of a compact operator on a Hilbert space.

\begin{Theorem}{Fredholm alternative theorem}
\label{thm:fredholm-alternative-compact}
Let $A\in\mathcal{L}(H)$ be compact. Then $\sigma(A)$ is a compact set having no limit point except perhaps $\lambda=0$. Furthermore, given any $\lambda\in\mathbb{C}\backslash\{0\}$, either\begin{enumerate}\item $\lambda\in\rho(A)$, or\item $\lambda\in\sigma_p(A)$ is an eigenvalue of finite multiplicity.\end{enumerate}
\end{Theorem}

\begingroup
\postdisplaypenalty=10000
\renewcommand{\endProof}{\par\nobreak\hfill$\blacksquare$\par\vspace{0.1cm}}
\begin{Proof}
Let $G=\mathbb{C}\backslash\{0\}$ and
\begin{equation}\label{eq:compact-fredholm-scaled-operator-family}
B(\lambda)=\dfrac{1}{\lambda}A.
\end{equation}
Then note that
\begin{equation}\label{eq:compact-fredholm-scaled-inverse-identity}
(\lambda I-A)^{-1}=\dfrac{1}{\lambda}\left(I-\dfrac{1}{\lambda}A\right)^{-1}=\dfrac{1}{\lambda}(I-B(\lambda))^{-1}.
\end{equation}
The result follows directly from Theorem~\ref{lem:analytic-finite-rank-alternative}.
\end{Proof}
\endgroup
